\documentclass{article}

\usepackage[margin=1in]{geometry}
\usepackage{graphicx, amsmath, amsthm, amssymb, enumitem, lipsum, hyperref}

\setlist[enumerate]{leftmargin=1.25cm}

\newcommand{\N}{\mathbb{N}}
\newcommand{\Z}{\mathbb{Z}}
\newcommand{\Q}{\mathbb{Q}}
\newcommand{\R}{\mathbb{R}}
\newcommand{\C}{\mathbb{C}}

\newcommand{\rotvert}{\rotatebox[origin=c]{90}{$\vert$}}

\newenvironment{note}[2][Note]{\begin{trivlist}
\item[\hskip \labelsep {\bfseries #1}\hskip \labelsep {\bfseries #2.}]}{\end{trivlist}}

\newenvironment{fact}[2][Fact]{\begin{trivlist}
\item[\hskip \labelsep {\bfseries #1}\hskip \labelsep {\bfseries #2.}]}{\end{trivlist}}

\newenvironment{definition}[2][Definition]{\begin{trivlist}
\item[\hskip \labelsep {\bfseries #1}\hskip \labelsep {\bfseries #2.}]}{\end{trivlist}}

\newenvironment{proposition}[2][Proposition]{\begin{trivlist}
\item[\hskip \labelsep {\bfseries #1}\hskip \labelsep {\bfseries #2.}]}{\end{trivlist}}

\newenvironment{principle}[2][Principle]{\begin{trivlist}
\item[\hskip \labelsep {\bfseries #1}\hskip \labelsep {\bfseries #2.}]}{\end{trivlist}}

\newenvironment{lemma}[2][Lemma]{\begin{trivlist}
\item[\hskip \labelsep {\bfseries #1}\hskip \labelsep {\bfseries #2.}]}{\end{trivlist}}

\newenvironment{theorem}[2][Theorem]{\begin{trivlist}
\item[\hskip \labelsep {\bfseries #1}\hskip \labelsep {\bfseries #2.}]}{\end{trivlist}}

\newenvironment{corollary}[2][Corollary]{\begin{trivlist}
\item[\hskip \labelsep {\bfseries #1}\hskip \labelsep {\bfseries #2.}]}{\end{trivlist}}

\newenvironment{envsection}[1]{\begin{trivlist}
\item[\hskip \labelsep {\bfseries #1}]}{\end{trivlist}}


\begin{document}
\setlength{\abovedisplayskip}{4pt}
\setlength{\belowdisplayskip}{4pt}


\begin{center}
    \Large{\textbf{Homework 4}}\\
    \normalsize February 19 - 26
\end{center}
\vspace{10pt}

\subsection*{Reading}

% \begin{envsection}{1. Rogers and Girolami:}
%     Sections 2.1--2.5. Most should be review, but I'd still like you to go through and make sure you understand it all. You can skip 2.5.2 and 2.5.4 if you'd like.
% \end{envsection}

% \begin{envsection}{1. Probabilistic Machine Learning}
%     Sections 2.1 - 2.2. You can skip 2.2.2.3, 2.2.5.3, 2.2.5.4, and 2.2.6.

%     2.6.1 and 2.6.2
% \end{envsection}

\begin{envsection}{1. CS229 Probability Reference:}
    Sections 1--4.1. Some of this may be review, but I'd still like you to go through and make sure you understand it. Their treatment of continuous random variables is a bit sparse/terse, so if you end up finding yourself lost at all, take a look at R+G 2.4, which is a bit gentler of an introduction.
\end{envsection}

\begin{envsection}{2. Rogers and Girolami:}
    Sections 2.2--2.5, excluding 2.5.2 and 2.5.4. I want you to \textit{very lightly} skim these sections after reading the CS229 notes, making sure you understand it all. You can also take a look at these sections if you get stuck on the CS229 notes for a gentler introduction to some topics.
\end{envsection}

\begin{envsection}{2. Rogers and Girolami:}
    Sections 2.1, 2.6--2.8.
\end{envsection}

\begin{envsection}{3. CS229 Notes:}
    If you have the time, go back and finish Section 8.1 (you can skip 8.1.1 if you want, but it is pretty interesting).
\end{envsection}

\subsection*{Problems}

\begin{envsection}{Problem 1.}
    
\end{envsection}

\subsection*{Fun Stuff}

\begin{envsection}{1. Watch:}
    \href{https://www.youtube.com/watch?v=D8GOeCFFby4&t=1593s}{``The most complex model we actually understand''} - Welch Labs
\end{envsection}


\end{document}